% GENERATED FILE. Edit canonical lessons, then run npm run generate:books.
% canonical-source-hash: fnv1a64:de7cb9e2a7b28c68
% canonical-lessons: TA-C368-kalviyarivu, TA-C368-vilippunarvu, TA-C368-utavittokai, TA-C368-karuttaranku, TA-C368-vina

\chapter{Learning for Everyone}
\label{ch:ta-learning-for-everyone}
\hlchaptermodality{368}

\begin{chapteropening}
\textbf{By the end of this chapter:} \emph{I can talk about literacy, awareness and scholarships.}

Complete the last lesson of chapter 368: I can talk about literacy, awareness and scholarships.
\end{chapteropening}

\section[\texorpdfstring{\ta{கல்வியறிவு} (kalviyaṟivu)}{கல்வியறிவு (kalviyaṟivu)}]{\ta{கல்வியறிவு} (kalviyaṟivu) — literacy}

\label{lesson:TA-C368-kalviyarivu}



\begin{quote}
Before the new one: say the Tamil for a genius, then the Tamil for to teach formally.
\end{quote}

\subsection*{You'll want to know: \ta{கல்வியறிவு}}
\textbf{\ta{கல்வியறிவு}} — \emph{kalviyaṟivu} — \textquotedblleft{}literacy\textquotedblright{}.

\begin{quote}
\ta{இந்த} \ta{மாவட்டத்தில்} \ta{கல்வியறிவு} \ta{அதிகரித்திருக்கிறது}.
\end{quote}

\begin{quote}
\emph{inta māvaṭṭattil kalviyaṟivu atikarittirukkiṟatu.}
\end{quote}

\begin{quote}
Literacy has risen in this district.
\end{quote}

\ta{கல்வி}, education + \ta{அறிவு}, knowledge.

\begin{cousinweb}
\textbf{\ta{கல்வியறிவு}} (\emph{kalviyaṟivu}) is native: \emph{kalvi}, \textquotedblleft{}learning, education\textquotedblright{}, from \emph{kal}, \textquotedblleft{}to learn\textquotedblright{} (DEDR 1297), and \emph{aṟivu}, \textquotedblleft{}knowledge\textquotedblright{} (DEDR 314).
\end{cousinweb}

\begin{grammarlens}[title={What you've built}]
One more word for giving reasons and opinions.
\end{grammarlens}

\subsection*{Your turn}
\begin{itemize}
\raggedright
  \item \emph{Say it:} \emph{kalviyaṟivu}
  \item \emph{Say it:} \emph{kalviyaṟivu}, once more
  \item \emph{Say it:} \emph{inta māvaṭṭattil kalviyaṟivu atikarittirukkiṟatu.} — Literacy has risen in this district.
  \item \emph{Recall:} say the Tamil for a genius, then the Tamil for to teach formally, then say \emph{kalviyaṟivu} again
\end{itemize}

\begin{tcolorbox}[breakable,colback=teal!4,colframe=teal!35!black,title={Before you move on}]
What does \ta{கல்வியறிவு} mean? (\textquotedblleft{}Literacy\textquotedblright{}.) Say it once more.
\end{tcolorbox}
\section[\texorpdfstring{\ta{விழிப்புணர்வு} (viḻippuṇarvu)}{விழிப்புணர்வு (viḻippuṇarvu)}]{\ta{விழிப்புணர்வு} (viḻippuṇarvu) — awareness}

\label{lesson:TA-C368-vilippunarvu}



\begin{quote}
Before the new one: say the Tamil for to teach formally, then the Tamil for literacy.
\end{quote}

\subsection*{You'll want to know: \ta{விழிப்புணர்வு}}
\textbf{\ta{விழிப்புணர்வு}} — \emph{viḻippuṇarvu} — \textquotedblleft{}awareness\textquotedblright{}.

\begin{quote}
\ta{சுத்தத்தைப்} \ta{பற்றி} \ta{மக்களிடையே} \ta{விழிப்புணர்வு} \ta{தேவை}.
\end{quote}

\begin{quote}
\emph{cuttattaip paṟṟi makkaḷiṭaiyē viḻippuṇarvu tēvai.}
\end{quote}

\begin{quote}
Awareness of cleanliness is needed among people.
\end{quote}

\ta{விழிப்புணர்வு} \ta{ஏற்படுத்து}, to raise awareness.

\begin{cousinweb}
\textbf{\ta{விழிப்புணர்வு}} (\emph{viḻippuṇarvu}) is native: \emph{viḻippu}, \textquotedblleft{}wakefulness\textquotedblright{}, from \emph{viḻi}, \textquotedblleft{}to open the eyes, wake\textquotedblright{} (DEDR 5429), and \emph{uṇarvu}, \textquotedblleft{}feeling, consciousness\textquotedblright{} (DEDR 603).
\end{cousinweb}

\begin{grammarlens}[title={What you've built}]
One more word for giving reasons and opinions.
\end{grammarlens}

\subsection*{Your turn}
\begin{itemize}
\raggedright
  \item \emph{Say it:} \emph{viḻippuṇarvu}
  \item \emph{Say it:} \emph{viḻippuṇarvu}, once more
  \item \emph{Say it:} \emph{cuttattaip paṟṟi makkaḷiṭaiyē viḻippuṇarvu tēvai.} — Awareness of cleanliness is needed among people.
  \item \emph{Recall:} say the Tamil for to teach formally, then the Tamil for literacy, then say \emph{viḻippuṇarvu} again
\end{itemize}

\begin{tcolorbox}[breakable,colback=teal!4,colframe=teal!35!black,title={Before you move on}]
What does \ta{விழிப்புணர்வு} mean? (\textquotedblleft{}Awareness\textquotedblright{}.) Say it once more.
\end{tcolorbox}
\section[\texorpdfstring{\ta{உதவித்தொகை} (utavittokai)}{உதவித்தொகை (utavittokai)}]{\ta{உதவித்தொகை} (utavittokai) — a scholarship; a grant}

\label{lesson:TA-C368-utavittokai}



\begin{quote}
Before the new one: say the Tamil for literacy, then the Tamil for awareness.
\end{quote}

\subsection*{You'll want to know: \ta{உதவித்தொகை}}
\textbf{\ta{உதவித்தொகை}} — \emph{utavittokai} — \textquotedblleft{}a scholarship; a grant\textquotedblright{}.

\begin{quote}
\ta{அவளுக்குப்} \ta{படிப்பதற்கு} \ta{உதவித்தொகை} \ta{கிடைத்தது}.
\end{quote}

\begin{quote}
\emph{avaḷukkup paṭippataṟku utavittokai kiṭaittatu.}
\end{quote}

\begin{quote}
She got a scholarship for her studies.
\end{quote}

\ta{உதவி}, help, + a word for a sum of money.

\begin{cousinweb}
\textbf{\ta{உதவித்தொகை}} (\emph{utavittokai}) is native: \emph{utavi}, \textquotedblleft{}help, aid\textquotedblright{}, from \emph{utavu} (DEDR 609), and \emph{tokai}, \textquotedblleft{}collection, sum, total\textquotedblright{} (DEDR 3476).
\end{cousinweb}

\begin{grammarlens}[title={What you've built}]
One more word for giving reasons and opinions.
\end{grammarlens}

\subsection*{Your turn}
\begin{itemize}
\raggedright
  \item \emph{Say it:} \emph{utavittokai}
  \item \emph{Say it:} \emph{utavittokai}, once more
  \item \emph{Say it:} \emph{avaḷukkup paṭippataṟku utavittokai kiṭaittatu.} — She got a scholarship for her studies.
  \item \emph{Recall:} say the Tamil for literacy, then the Tamil for awareness, then say \emph{utavittokai} again
\end{itemize}

\begin{tcolorbox}[breakable,colback=teal!4,colframe=teal!35!black,title={Before you move on}]
What does \ta{உதவித்தொகை} mean? (\textquotedblleft{}A scholarship; a grant\textquotedblright{}.) Say it once more.
\end{tcolorbox}
\section[\texorpdfstring{\ta{கருத்தரங்கு} (karuttaraṅku)}{கருத்தரங்கு (karuttaraṅku)}]{\ta{கருத்தரங்கு} (karuttaraṅku) — a seminar}

\label{lesson:TA-C368-karuttaranku}



\begin{quote}
Before the new one: say the Tamil for awareness, then the Tamil for a scholarship.
\end{quote}

\subsection*{You'll want to know: \ta{கருத்தரங்கு}}
\textbf{\ta{கருத்தரங்கு}} — \emph{karuttaraṅku} — \textquotedblleft{}a seminar\textquotedblright{}.

\begin{quote}
\ta{நாளை} \ta{பல்கலைக்கழகத்தில்} \ta{ஒரு} \ta{கருத்தரங்கு} \ta{நடக்கிறது}.
\end{quote}

\begin{quote}
\emph{nāḷai palkalaikkaḻakattil oru karuttaraṅku naṭakkiṟatu.}
\end{quote}

\begin{quote}
A seminar is being held at the university tomorrow.
\end{quote}

\ta{கருத்து}, opinion + \ta{அரங்கு}, hall: a hall for ideas.

\begin{cousinweb}
\textbf{\ta{கருத்தரங்கு}} (\emph{karuttaraṅku}) is a modern coinage: native \emph{karuttu}, \textquotedblleft{}idea, opinion\textquotedblright{}, from \emph{karutu}, \textquotedblleft{}to consider\textquotedblright{} (DEDR 1283), and \emph{araṅku}, \textquotedblleft{}stage, assembly hall\textquotedblright{}, which DEDR lists among Indo-Aryan loans (DEDR App. 8).
\end{cousinweb}

\begin{grammarlens}[title={What you've built}]
One more word for giving reasons and opinions.
\end{grammarlens}

\subsection*{Your turn}
\begin{itemize}
\raggedright
  \item \emph{Say it:} \emph{karuttaraṅku}
  \item \emph{Say it:} \emph{karuttaraṅku}, once more
  \item \emph{Say it:} \emph{nāḷai palkalaikkaḻakattil oru karuttaraṅku naṭakkiṟatu.} — A seminar is being held at the university tomorrow.
  \item \emph{Recall:} say the Tamil for awareness, then the Tamil for a scholarship, then say \emph{karuttaraṅku} again
\end{itemize}

\begin{tcolorbox}[breakable,colback=teal!4,colframe=teal!35!black,title={Before you move on}]
What does \ta{கருத்தரங்கு} mean? (\textquotedblleft{}A seminar\textquotedblright{}.) Say it once more.
\end{tcolorbox}
\section[\texorpdfstring{\ta{வினா} (viṉā)}{வினா (viṉā)}]{\ta{வினா} (viṉā) — a question}

\label{lesson:TA-C368-vina}



\begin{quote}
Before the new one: say the Tamil for a scholarship, then the Tamil for a seminar.
\end{quote}

\subsection*{You'll want to know: \ta{வினா}}
\textbf{\ta{வினா}} — \emph{viṉā} — \textquotedblleft{}a question\textquotedblright{}.

\begin{quote}
\ta{தேர்வில்} \ta{பத்து} \ta{வினாக்கள்} \ta{இருந்தன}.
\end{quote}

\begin{quote}
\emph{tērvil pattu viṉākkaḷ iruntaṉa.}
\end{quote}

\begin{quote}
There were ten questions in the exam.
\end{quote}

Plural: \ta{வினாக்கள்}, \emph{viṉākkaḷ}.

A written word for \ta{கேள்வி}; \ta{வினாத்தாள்}, a question paper.

\begin{cousinweb}
\textbf{\ta{வினா}} (\emph{viṉā}) is native Dravidian: DEDR 5516 lists Tamil \emph{viṉā}, \textquotedblleft{}to ask, investigate, listen; question\textquotedblright{}, with Malayalam \emph{vinavuka}, \textquotedblleft{}to ask\textquotedblright{}.
\end{cousinweb}

\begin{grammarlens}[title={What you've built}]
One more word for giving reasons and opinions.
\end{grammarlens}

\subsection*{Your turn}
\begin{itemize}
\raggedright
  \item \emph{Say it:} \emph{viṉā}
  \item \emph{Say it:} \emph{viṉā}, once more
  \item \emph{Say it:} \emph{tērvil pattu viṉākkaḷ iruntaṉa.} — There were ten questions in the exam.
  \item \emph{Recall:} say the Tamil for a scholarship, then the Tamil for a seminar, then say \emph{viṉā} again
\end{itemize}

\begin{tcolorbox}[breakable,colback=teal!4,colframe=teal!35!black,title={Before you move on}]
What does \ta{வினா} mean? (\textquotedblleft{}A question\textquotedblright{}.) Say it once more.
\end{tcolorbox}
